% ENEM 2019-2023 — 9 questões MCQ — Funções I
% Fonte: lista "Funções I" do site projetoagathaedu.com.br
% Omitidas: Q3,Q5,Q6,Q7,Q10,Q11,Q14,Q15,Q16 (imagens)
%           Q19-Q22 (HTML truncado — não processadas)

---
tipo: Múltipla Escolha
dificuldade: Média
nivel: médio
disciplina: Matemática
assunto: Função Quadrática
gabarito: D
tags: [empresa, desempenho, vendas, mínimo, parábola]
fonte: concurso
concurso: ENEM
banca: INEP
ano: 2023
numero: 1
---
\question
Analisando as vendas de uma empresa, o gerente concluiu que o montante arrecadado, em milhar de real, poderia ser calculado pela expressão \( V(x) = \frac{x^2}{4} - 10x + 105 \), em que os valores de \( x \) representam os dias do mês, variando de 1 a 30.

Um dos fatores para avaliar o desempenho mensal da empresa é verificar qual é o menor montante diário \( V_0 \) arrecadado ao longo do mês e classificar o desempenho conforme as categorias apresentadas a seguir, em que as quantidades estão expressas em milhar de real.

-- Ótimo: \( V_0 \geq 24 \)

-- Bom: \( 20 \leq V_0 < 24 \)

-- Normal: \( 10 \leq V_0 < 20 \)

-- Ruim: \( 4 \leq V_0 < 10 \)

-- Péssimo: \( V_0 < 4 \)

No caso analisado, qual seria a classificação do desempenho da empresa?
\begin{choices}
\choice Ótimo.
\choice Bom.
\choice Normal.
\correctchoice Ruim.
\choice Péssimo.
\end{choices}

---
tipo: Múltipla Escolha
dificuldade: Fácil
nivel: médio
disciplina: Matemática
assunto: Função Afim
gabarito: D
tags: [concreto, laje, bombeamento, taxa fixa]
fonte: concurso
concurso: ENEM
banca: INEP
ano: 2023
numero: 2
---
\question
Para concretar a laje de sua residência, uma pessoa contratou uma construtora. Tal empresa informa que o preço \( y \) do concreto bombeado é composto de duas partes: uma fixa, chamada de taxa de bombeamento, e outra variável, que depende do volume \( x \) de concreto utilizado. Sabe-se que a taxa de bombeamento custa R\$ 500,00 e que o metro cúbico do concreto bombeado é de R\$ 250,00.

A expressão que representa o preço \( y \) em função do volume \( x \), em metro cúbico, é
\begin{oneparchoices}
\choice \( y = 250x \)
\choice \( y = 500x \)
\choice \( y = 750x \)
\correctchoice \( y = 250x + 500 \)
\choice \( y = 500x + 250 \)
\end{oneparchoices}

---
tipo: Múltipla Escolha
dificuldade: Fácil
nivel: médio
disciplina: Matemática
assunto: Função Afim
gabarito: B
tags: [torneira, balde, gotas, notação científica]
fonte: concurso
concurso: ENEM
banca: INEP
ano: 2020
numero: 4
---
\question
Uma torneira está gotejando água em um balde com capacidade de 18 litros. No instante atual, o balde se encontra com ocupação de 50\% de sua capacidade. A cada segundo caem 5 gotas de água da torneira, e uma gota é formada, em média, por \( 5 \times 10^{-2} \) mL de água.

Quanto tempo, em hora, será necessário para encher completamente o balde, partindo do instante atual?
\begin{oneparchoices}
\choice \( 2 \times 10^{1} \)
\correctchoice \( 1 \times 10^{1} \)
\choice \( 2 \times 10^{-2} \)
\choice \( 1 \times 10^{-2} \)
\choice \( 1 \times 10^{-3} \)
\end{oneparchoices}

---
tipo: Múltipla Escolha
dificuldade: Fácil
nivel: médio
disciplina: Matemática
assunto: Função Afim
gabarito: B
tags: [soja, produtor rural, hectare, lucro]
fonte: concurso
concurso: ENEM
banca: INEP
ano: 2020
numero: 8
---
\question
Por muitos anos, o Brasil tem figurado no cenário mundial entre os maiores produtores e exportadores de soja. Entre os anos de 2010 e 2014, houve uma forte tendência de aumento da produtividade, porém, um aspecto dificultou esse avanço: o alto custo do imposto ao produtor associado ao baixo preço de venda do produto. Em média, um produtor gastava R\$ 1.200,00 por hectare plantado, e vendia por R\$ 50,00 cada saca de 60 kg. Ciente desses valores, um produtor pode, em certo ano, determinar uma relação do lucro \( L \) que obteve em função das sacas de 60 kg vendidas. Suponha que ele plantou 10 hectares de soja em sua propriedade, na qual colheu \( x \) sacas de 60 kg e todas as sacas foram vendidas.

Qual é a expressão que determinou o lucro \( L \) em função de \( x \) obtido por esse produtor nesse ano?
\begin{oneparchoices}
\choice \( L(x) = 50x - 1{.}200 \)
\correctchoice \( L(x) = 50x - 12{.}000 \)
\choice \( L(x) = 50x + 12{.}000 \)
\choice \( L(x) = 500x - 1{.}200 \)
\choice \( L(x) = 1{.}200x - 500 \)
\end{oneparchoices}

---
tipo: Múltipla Escolha
dificuldade: Média
nivel: médio
disciplina: Matemática
assunto: Função Quadrática
gabarito: D
tags: [chocolate, barras, lucro, vértice, máximo]
fonte: concurso
concurso: ENEM
banca: INEP
ano: 2020
numero: 9
---
\question
Uma empresa de chocolates consultou o gerente de produção e verificou que existem cinco tipos diferentes de barras de chocolate que podem ser produzidas, com os seguintes preços no mercado:

-- Barra I: R\$ 2,00;

-- Barra II: R\$ 3,50;

-- Barra III: R\$ 4,00;

-- Barra IV: R\$ 7,00;

-- Barra V: R\$ 8,00.

Analisando as tendências do mercado, que incluem a quantidade vendida e a procura pelos consumidores, o gerente de vendas da empresa verificou que o lucro \( L \) com a venda de barras de chocolate é expresso pela função \( L(x) = -x^2 + 14x - 45 \), em que \( x \) representa o preço da barra de chocolate.

A empresa decide investir na fabricação da barra de chocolate cujo preço praticado no mercado renderá o maior lucro.

Nessas condições, a empresa deverá investir na produção da barra
\begin{oneparchoices}
\choice I.
\choice II.
\choice III.
\correctchoice IV.
\choice V.
\end{oneparchoices}

---
tipo: Múltipla Escolha
dificuldade: Média
nivel: médio
disciplina: Matemática
assunto: Função Exponencial
gabarito: B
tags: [bactérias, reprodução, duplicação, laboratório]
fonte: concurso
concurso: ENEM-PPL
banca: INEP
ano: 2020
numero: 12
---
\question
Um laboratório realizou um teste para calcular a velocidade de reprodução de um tipo de bactéria. Para tanto, realizou um experimento para observar a reprodução de uma quantidade \( x \) dessas bactérias por um período de duas horas. Após esse período, constava no habitáculo do experimento uma população de 189.440 da citada bactéria. Constatou-se, assim, que a população de bactérias dobrava a cada 0,25 hora.

A quantidade inicial de bactérias era de
\begin{oneparchoices}
\choice 370.
\correctchoice 740.
\choice 1.480.
\choice 11.840.
\choice 23.680.
\end{oneparchoices}

---
tipo: Múltipla Escolha
dificuldade: Média
nivel: médio
disciplina: Matemática
assunto: Função Exponencial
gabarito: E
tags: [PIS, Cofins, imposto cumulativo, circulação de mercadorias]
fonte: concurso
concurso: ENEM-PPL
banca: INEP
ano: 2020
numero: 13
---
\question
Um imposto é dito cumulativo se incide em duas ou mais etapas da circulação de mercadorias, sem que na etapa posterior possa ser abatido o montante pago na etapa anterior. PIS e Cofins são exemplos de impostos cumulativos e correspondem a um percentual total de 3,65\%, que incide em cada etapa da comercialização de um produto.

Considere um produto com preço inicial \( C \). Suponha que ele é revendido para uma loja pelo preço inicial acrescido dos impostos descritos. Em seguida, o produto é revendido por essa loja ao consumidor pelo valor pago acrescido novamente dos mesmos impostos.

Qual a expressão algébrica que corresponde ao valor pago em impostos pelo consumidor?
\begin{choices}
\choice \( C \times 0{,}0365 \)
\choice \( 2C \times 0{,}0365 \)
\choice \( C \times 1{,}0365^2 \)
\choice \( C \times (1 + 2 \times 0{,}0365) \)
\correctchoice \( 2C \times 0{,}0365 + C \times 0{,}0365^2 \)
\end{choices}

---
tipo: Múltipla Escolha
dificuldade: Média
nivel: médio
disciplina: Matemática
assunto: Função Logarítmica
gabarito: D
tags: [jardineiro, planta ornamental, crescimento, logaritmo]
fonte: concurso
concurso: ENEM-PPL
banca: INEP
ano: 2019
numero: 17
---
\question
Um jardineiro cultiva plantas ornamentais e as coloca à venda quando estas atingem 30 centímetros de altura. Esse jardineiro estudou o crescimento de suas plantas, em função do tempo, e deduziu uma fórmula que calcula a altura em função do tempo, a partir do momento em que a planta brota do solo até o momento em que ela atinge sua altura máxima de 40 centímetros. A fórmula é \( h = 5 \cdot \log_2(t + 1) \), em que \( t \) é o tempo contado em dia e \( h \), a altura da planta em centímetro.

A partir do momento em que uma dessas plantas é colocada à venda, em quanto tempo, em dia, ela alcançará sua altura máxima?
\begin{oneparchoices}
\choice 63
\choice 96
\choice 128
\correctchoice 192
\choice 255
\end{oneparchoices}

---
tipo: Múltipla Escolha
dificuldade: Difícil
nivel: médio
disciplina: Matemática
assunto: Funções
gabarito: A
tags: [mola, constante elástica, espiral, módulo de elasticidade]
fonte: concurso
concurso: ENEM-PPL
banca: INEP
ano: 2019
numero: 18
---
\question
Para certas molas, a constante elástica (\( C \)) depende do diâmetro médio da circunferência da mola (\( D \)), do número de espirais úteis (\( N \)), do diâmetro (\( d \)) do fio de metal do qual é formada a mola e do módulo de elasticidade do material (\( G \)). A fórmula evidencia essas relações de dependência:
\[ C = \frac{G \cdot d^4}{8 \cdot D^3 \cdot N} \]

O dono de uma fábrica possui uma mola \( M_1 \) em um de seus equipamentos, que tem características \( D_1 \), \( d_1 \), \( N_1 \) e \( G_1 \), com uma constante elástica \( C_1 \). Essa mola precisa ser substituída por outra, \( M_2 \), produzida com outro material e com características diferentes, bem como uma nova constante elástica \( C_2 \), da seguinte maneira:

I) \( D_2 = \dfrac{D_1}{3} \);

II) \( d_2 = 3d_1 \);

III) \( N_2 = 9N_1 \).

Além disso, a constante de elasticidade \( G_2 \) do novo material é igual a \( 4G_1 \).

O valor da constante \( C_2 \) em função da constante \( C_1 \) é
\begin{choices}
\correctchoice \( C_2 = 972 \cdot C_1 \)
\choice \( C_2 = 108 \cdot C_1 \)
\choice \( C_2 = 4 \cdot C_1 \)
\choice \( C_2 = \dfrac{4}{3} \cdot C_1 \)
\choice \( C_2 = \dfrac{4}{9} \cdot C_1 \)
\end{choices}
